% =========================================================================
%  Capítulo 4 — Conclusão
% =========================================================================

\chapter{Conclusão}
\label{cap:conclusao}

% --- INSTRUÇÕES (apagar antes da entrega) ---------------------
% A conclusão NÃO é um resumo do que se fez. É uma reflexão
% crítica sobre:
%   - O que ficou bem feito.
%   - O que ficaria diferente se tivessem mais tempo.
%   - Que aprendizagens foram feitas.
%   - Que trabalho futuro faz sentido.
% Evitar frases vazias do tipo ``foi um projeto enriquecedor''.
% Concretizar: o que aprenderam SOBRE arquitetura web que não
% sabiam antes?
% ---------------------------------------------------------------


\section{Síntese do trabalho realizado}
\label{sec:sintese}

% Síntese curta do que foi efetivamente entregue. Esta secção
% deve ser objetiva: percentagem de requisitos implementados,
% módulos funcionais, etc.

[Apresentar de forma sintética o que foi efetivamente
implementado. Recomenda-se referir a cobertura dos requisitos
funcionais (RF-01 a RF-15) e dos requisitos não funcionais,
indicando claramente o que ficou completo, parcial e por
implementar.]


\section{Avaliação dos critérios de aceitação}
\label{sec:criterios}

% Confrontar o trabalho realizado com os critérios de aceitação
% do enunciado. Esta tabela facilita a vida ao avaliador.

A Tabela~\ref{tab:criterios} confronta o trabalho realizado com
os critérios de aceitação definidos no enunciado.

\begin{table}[H]
\centering
\caption{Verificação dos critérios de aceitação do projeto.}
\label{tab:criterios}
\footnotesize
\renewcommand{\arraystretch}{1.4}
\begin{tabularx}{\linewidth}{l X c}
\toprule
\textbf{ID} & \textbf{Critério} & \textbf{Estado} \\
\midrule
CA-01 & Um utilizador autenticado acede apenas às operações permitidas pelo seu perfil. & \\
CA-02 & O sistema suporta os três tipos de plano de cultivo e a sua associação a lotes. & \\
CA-03 & O registo de medições ambientais pode originar alertas automáticos. & \\
CA-04 & O sistema demonstra comportamento distinto em modo \emph{Manual} e \emph{Automático}. & \\
CA-05 & É possível consultar histórico, exportar relatórios e observar rastreabilidade. & \\
CA-06 & A arquitetura de armazenamento no \emph{browser} está implementada e coerente. & \\
CA-07 & A implementação está alinhada com a especificação OpenAPI. & \\
\bottomrule
\end{tabularx}
\end{table}

\noindent
\textit{Legenda:} preencher cada linha com \checkmark\ (cumprido),
$\sim$ (parcialmente cumprido) ou $\times$ (não cumprido), seguido
de uma breve evidência.


\section{Reflexão crítica}
\label{sec:reflexao}

% Esta é a secção que distingue um relatório bom de um relatório
% medíocre. Ser HONESTO. Os avaliadores valorizam mais um
% relatório que reconhece limitações do que um que pretende ser
% perfeito.

\subsection{O que correu bem}

[Identificar os aspetos do projeto onde a equipa considera ter
tomado boas decisões. Justificar com evidências concretas — não
basta dizer ``a arquitetura ficou bem''; é preciso explicar
porquê.]

\subsection{O que faríamos diferente}

[Identificar honestamente as decisões que, em retrospetiva, a
equipa repensaria. Esta é uma das partes mais valorizadas do
relatório: mostra capacidade de aprendizagem e pensamento
crítico.]

\subsection{Aprendizagens da equipa}

[Aprendizagens técnicas e de processo. Que conceitos ficaram mais
claros? Que ferramentas mostraram ser mais ou menos adequadas? Que
prática de equipa funcionou melhor?]


\section{Trabalho futuro}
\label{sec:futuro}

% Lista priorizada de melhorias possíveis. Pode incluir:
%   - Requisitos identificados mas não implementados.
%   - Otimizações de desempenho ou usabilidade.
%   - Funcionalidades que extravasam o âmbito atual mas
%     fariam sentido numa próxima iteração.

Identificam-se como direções relevantes para iterações futuras:

\begin{itemize}
  \item {}[Listar funcionalidades não implementadas que fariam sentido — e.g., notificações \emph{push}, \emph{dashboards} avançados, integração com sensores reais.]
  \item {}[Listar melhorias de natureza técnica — e.g., \emph{cobertura de testes automatizados, observabilidade, hardening de segurança}.]
  \item {}[Listar mudanças de processo ou metodologia que ajudariam num projeto semelhante.]
\end{itemize}
